%% ============================================================
%%  CAPITOLO 1 — Il contesto clinico delle Note AIFA
%%  Tesi triennale — Francesco Zamar, UniTS, a.a. 2025/2026
%%  Relatore: Prof. Andrea De Lorenzo
%%  Stile: prosa narrativa fluente, nessun itemize/enumerate.
%%  Macro canoniche da macros_finale.tex.
%% ============================================================

\chapter{Il contesto clinico delle Note AIFA}
\label{ch:contesto}

Il Servizio Sanitario Nazionale (SSN) italiano garantisce l'accesso ai farmaci essenziali
attraverso un sistema di rimborsabilità che tiene insieme sostenibilità economica e
appropriatezza clinica.
Per un sottoinsieme rilevante dei medicinali di fascia~A, la rimborsabilità non è
automatica: il prescrittore deve verificare che il paziente soddisfi i criteri
definiti in atti normativi di dettaglio denominati \emph{Note AIFA}.
Questa tesi si occupa di automatizzare tale verifica per quattro Note selezionate,
costruendo un sistema di supporto alla decisione clinica capace di produrre un
esito motivato e tracciabile a partire dal profilo clinico del paziente e dal testo
normativo ufficiale~\cite{aifa2024note}.
Il corpus di valutazione comprende \goldn\ casi gold distribuiti su 122 scenari
clinici realistici; la distribuzione delle etichette riflette l'eterogeneità
delle situazioni prescrittive: \goldRimb\ casi \decision{RIMBORSABILE},
\goldNonRimb\ casi \decision{NON\_RIMBORSABILE}, \goldNonDet\ casi
\decision{NON\_DETERMINABILE} e \goldRouted\ casi instradati verso una Nota
dipendente.

% ------------------------------------------------------------------
\section{Sistemi di supporto alla decisione clinica}
\label{sec:cdss-overview}
% ------------------------------------------------------------------

Un sistema di supporto alla decisione clinica (CDSS, \emph{Clinical Decision
Support System}) è un software progettato per assistere il professionista sanitario
nell'elaborazione di informazioni cliniche al fine di migliorare la qualità
delle decisioni~\cite{greenes2007cdss}.
La definizione è intenzionalmente ampia: include sistemi di allerta per le
interazioni farmacologiche, strumenti di scoring prognostico, motori di verifica
normativa e raccomandatori terapeutici.
La distinzione rilevante per questo lavoro è quella tra sistemi attivi, che
intervengono proattivamente nel flusso di lavoro clinico, e sistemi passivi, che
rispondono a una richiesta puntuale del medico~\cite{berner2007cdss}.

La storia dei CDSS inizia nel 1974 con MYCIN~\cite{shortliffe1974mycin}, un
sistema per la selezione della terapia antibiotica sviluppato a Stanford da
Shortliffe.
MYCIN codificava la conoscenza clinica esperta come insieme di regole di produzione
\emph{if-then} pesate con fattori di certezza, e dimostrava per la prima volta che
un software poteva raggiungere prestazioni paragonabili a quelle del clinico esperto
su un dominio ben delimitato.
Negli anni successivi, il modello si è evoluto verso linguaggi formali interoperabili:
l'Arden Syntax~\cite{hripcsak1994arden}, standardizzato da Health Level Seven (HL7) nel 1992, ha fornito
un framework per specificare moduli di logica clinica integrabili nelle cartelle
elettroniche, spostando l'enfasi dalla knowledge base monolitica alla modularità
delle regole.
Il decennio 2010 ha visto l'integrazione progressiva di metodi probabilistici
e di apprendimento automatico nei CDSS, con guadagni sulla generalizzazione ma
perdite sulla tracciabilità: un modello che impara pattern dai dati non produce
per costruzione una catena di giustificazioni verificabile~\cite{sutton2020overview}.

Il contesto delle Note AIFA pone requisiti specifici che limitano l'applicabilità
dei soli modelli probabilistici.
La decisione di rimborsabilità ha effetto economico e amministrativo diretto: ogni
prescrizione non conforme genera oneri sul SSN e può comportare responsabilità
disciplinare per il medico.
Una componente rule-based deterministica è dunque indispensabile per garantire
l'auditabilità del percorso decisionale.
Al tempo stesso, i dati clinici reali sono raramente completi, e la Nota può
contemplare parametri non rilevati al momento della prescrizione: è necessario
gestire formalmente l'incertezza senza rinunciare alla tracciabilità.
Il sistema descritto in questa tesi risponde a questi requisiti combinando un
motore di regole simbolico con un modulo di Retrieval-Augmented Generation, il
cui ruolo complementare è analizzato a partire dal \Cref{ch:progettazione}.

% ------------------------------------------------------------------
\section{La regolamentazione AIFA della rimborsabilità}
\label{sec:aifa-regolamentazione}
% ------------------------------------------------------------------

L'Agenzia Italiana del Farmaco (AIFA) è l'ente pubblico indipendente, posto
sotto la vigilanza del Ministero della Salute e del Ministero dell'Economia e delle
Finanze, che regola l'intero ciclo di vita del medicinale in Italia: valutazione
tecnico-scientifica, negoziazione del prezzo con il produttore, classificazione in
fascia di rimborsabilità e farmacovigilanza post-marketing~\cite{aifa2024note}.
La spesa farmaceutica pubblica si attesta intorno ai quindici miliardi di euro
annui, ed è il perimetro all'interno del quale l'AIFA esercita il suo mandato di
sostenibilità.

I medicinali commercializzati in Italia vengono classificati in tre grandi fasce:
la fascia~A comprende i farmaci essenziali e quelli per patologie croniche ad elevato
impatto sociale, erogati gratuitamente salvo eventuale ticket regionale; la fascia~H
riguarda i medicinali riservati all'uso ospedaliero o specialistico, anch'essi
a carico del SSN nelle strutture pubbliche; la fascia~C include i farmaci il cui
costo ricade interamente sul cittadino.
Per un sottoinsieme significativo dei farmaci di fascia~A, la rimborsabilità
non è incondizionata ma vincolata al rispetto dei criteri definiti nelle
\emph{Note AIFA}, atti pubblicati in Gazzetta Ufficiale e aggiornati con cadenza
variabile al ritmo delle evidenze cliniche.
Attualmente sono vigenti oltre cinquanta Note, numerate progressivamente, che coprono aree terapeutiche eterogenee:
cardiovascolare, metabolica, neurologica, reumatologica, gastroenterologica,
oncologica.

Ogni Nota è identificata da un numero progressivo e specifica, per una o più
categorie farmacologiche, le indicazioni cliniche rimborsabili, le controindicazioni
assolute, i percorsi terapeutici di primo e di secondo livello e i parametri di
stratificazione del rischio.
La complessità strutturale delle Note è molto variabile: una Nota semplice si
riduce a una singola condizione booleana (il farmaco è rimborsabile se la diagnosi
corrisponde a un codice ICD (\emph{International Classification of Diseases}) specifico), mentre una Nota complessa come la
\notenova\ richiede il calcolo di uno score composito, la verifica di soglie
differenziali per sesso, la presenza di una tabella di dosaggio con criteri di
riduzione multipla e la gestione di controindicazioni assolute che escludono intere
sotto-classi terapeutiche.
Tale eterogeneità strutturale è costitutiva, non accessoria: è esattamente il
motivo per cui la sola memorizzazione delle regole da parte del prescrittore è
una strategia insufficiente su scala.

% ------------------------------------------------------------------
\section{Le quattro Note oggetto di studio}
\label{sec:quattro-note}
% ------------------------------------------------------------------

Il sistema implementato in questa tesi codifica quattro Note AIFA selezionate
secondo un criterio di eterogeneità clinica e di complessità logica crescente.
La \notezero\ costituisce il caso logicamente più semplice, riducibile a una
congiunzione di due disgiunzioni.
La \notetredici\ introduce la stratificazione del rischio cardiovascolare con
soglie di colesterolo LDL (\emph{Low-Density Lipoprotein}) variabili per classe di rischio.
La \notesessanta\ aggiunge liste chiuse di farmaci e controindicazioni
farmaco-specifiche per singola molecola.
La \notenova\ è la più articolata dell'insieme: score composito, aritmetica
d'intervallo, soglie differenziali per sesso e tabella di dosaggio con criteri
di riduzione multipla~\cite{aifa2024note}.
Complessivamente, le quattro Note coprono 122 scenari gold con distribuzione
\decision{RIMB}/\decision{NON\_RIMB}/\decision{NON\_DET}/\decision{ROUTED}
pari a 69/39/9/5 casi.

% - - - - - - - - - - - - - - - - - - - - - - -
\subsection{Nota \notezero: gastroprotezione cronica}
\label{subsec:nota01}
% - - - - - - - - - - - - - - - - - - - - - - -

La \notezero\ regola la rimborsabilità degli inibitori di pompa protonica (IPP)
nella prevenzione delle complicanze gravi del tratto gastrointestinale superiore
in pazienti in terapia cronica con farmaci ulcerogeni~\cite{nota01aifa}.
Gli IPP agiscono bloccando irreversibilmente la pompa
H\textsuperscript{+}/K\textsuperscript{+}-ATPasi delle cellule parietali gastriche,
sopprimendo la secrezione acida basale e stimolata; i cinque principi attivi
rimborsabili ai sensi della Nota sono omeprazolo, esomeprazolo, lansoprazolo,
pantoprazolo e rabeprazolo.

Il razionale clinico della Nota poggia sul rischio di ospedalizzazione per
complicanze gravi (emorragia, perforazione, ostruzione pilorica), stimato tra
l'1 e il 2\% annuo nella popolazione generale in terapia con antinfiammatori non steroidei (FANS), e aumentato
fino a quattro-cinque volte nelle categorie a rischio specificate dalla Nota stessa.
La condizione necessaria di contesto è la co-somministrazione di un FANS o di
acido acetilsalicilico a basse dosi; la condizione sufficiente per la rimborsabilità
è la presenza di almeno uno tra i fattori di rischio elencati, fra i quali figurano
pregressa emorragia digestiva, ulcera peptica non guarita, terapia anticoagulante
concomitante, terapia corticosteroidea concomitante ed età avanzata (criterio
asserted dal clinico, non definito numericamente nel testo normativo).
La struttura logica è dunque una congiunzione tra la condizione di contesto e
la disgiunzione dei fattori di rischio: la semplicità di questa topologia non
elimina però la complessità operativa, poiché la Nota~\notezero\ presenta un
meccanismo di \emph{routing} inter-Nota, in cui la rimborsabilità dell'associazione
diclofenac/misoprostolo dipende dall'esito della \notesessanta\, introducendo
una dipendenza tra documenti normativi distinti.

% - - - - - - - - - - - - - - - - - - - - - - -
\subsection{Nota \notetredici: ipolipemizzanti nella dislipidemia}
\label{subsec:nota13}
% - - - - - - - - - - - - - - - - - - - - - - -

La \notetredici\ regola la rimborsabilità degli ipolipemizzanti nel trattamento
delle dislipidemie, introducendo un principio di proporzionalità tra rischio
cardiovascolare e soglia di trattamento~\cite{aifa2018nota13}.
L'ipercolesterolemia è il principale fattore di rischio cardiovascolare
modificabile su scala di popolazione: una riduzione dell'LDL-C di 1~mmol/L
si associa a una riduzione del rischio di eventi cardiovascolari maggiori di
circa il 22\%~\cite{atp2014cardio}.
La Nota non prescrive di trattare ogni paziente iperlipemico, ma afferma che
la soglia dell'LDL-C che giustifica la rimborsabilità è tanto più bassa quanto
più elevato è il rischio cardiovascolare stimato: un paziente a rischio molto
alto deve raggiungere un target LDL inferiore a 70~mg/dL, mentre un paziente
a rischio medio può restare non trattato farmacologicamente se il colesterolo
non supera 130~mg/dL.

I farmaci in scope comprendono le statine (simvastatina, pravastatina, fluvastatina,
lovastatina, atorvastatina, rosuvastatina), l'ezetimibe, i fibrati e gli acidi
grassi polinsaturi omega-3.
La stratificazione del rischio segue le linee guida della \emph{European Society of Cardiology} (ESC) e della \emph{European Atherosclerosis Society} (EAS)~\cite{escic2019dyslip}
e definisce cinque classi: basso (rischio SCORE $\leq$1\%), medio (2--3\%),
moderato (4--5\%), alto (5--10\%) e molto alto ($\geq$10\% o malattia
cardiovascolare conclamata), con target LDL rispettivamente di sola modifica
dello stile di vita, inferiore a 130, 115, 100 e 70~mg/dL.
Un'eccezione trasversale riguarda i pazienti con dislipidemie familiari
monogeniche: per questi, indipendentemente dallo score calcolato, i farmaci
di primo e secondo livello sono rimborsabili senza la necessità di dimostrare
il mancato raggiungimento del target con la sola dieta.
Questa regola di \emph{bypass} a priorità superiore rispecchia la logica clinica
secondo cui il rischio genetico assoluto giustifica l'intervento farmacologico
indipendentemente dal rischio epidemiologico calcolato.

% - - - - - - - - - - - - - - - - - - - - - - -
\subsection{Nota \notesessanta: antinfiammatori non steroidei}
\label{subsec:nota66}
% - - - - - - - - - - - - - - - - - - - - - - -

La \notesessanta\ regola la rimborsabilità degli antinfiammatori non steroidei
nella gestione del dolore e dell'infiammazione in quattro contesti
clinici specifici: artropatie su base connettivitica, osteoartrosi in fase
algica o infiammatoria, dolore neoplastico, attacco acuto di gotta~\cite{nota66aifa}.
I FANS agiscono mediante inibizione della ciclo-ossigenasi (COX) dell'acido
arachidonico, bloccando la sintesi di prostaglandine e trombossani responsabili
della risposta infiammatoria; la distinzione farmacologicamente rilevante è
tra gli inibitori non selettivi della COX, che espongono il paziente a rischio
gastrointestinale elevato, e i coxib (celecoxib, etoricoxib), inibitori
selettivi della COX-2 a minor rischio gastrointestinale ma con profilo di
rischio cardiovascolare aumentato.

La Nota disciplina 28 principi attivi, e la rimborsabilità non è un attributo
della categoria FANS in generale ma di ciascun principio attivo singolarmente.
La nimesulide è ammessa solo per il trattamento di breve durata del dolore
acuto; l'associazione ibuprofene/codeina è rimborsabile quando il dolore non
sia adeguatamente controllato con analgesici singoli; i coxib sono esclusi
dalla rimborsabilità nei pazienti con malattia cardiovascolare conclamata o
insufficienza renale grave, indipendentemente dall'indicazione.
Quest'ultimo meccanismo, una controindicazione farmaco-specifica che sovrascrive
l'esito positivo della verifica dell'indicazione, è la caratteristica logicamente
più impegnativa della Nota e motiva l'introduzione dell'operatore
\texttt{IN}$(\text{farmaco},\,\mathcal{L}_\text{FANS})$ nel linguaggio
del motore di regole.

% - - - - - - - - - - - - - - - - - - - - - - -
\subsection{Nota \notenova: anticoagulanti orali nella fibrillazione atriale}
\label{subsec:nota97}
% - - - - - - - - - - - - - - - - - - - - - - -

La \notenova\ regola la rimborsabilità di quattro anticoagulanti orali diretti
(DOAC), ossia dabigatran, rivaroxaban, apixaban ed edoxaban, nella fibrillazione
atriale non valvolare (FANV)~\cite{aifa2018nota97}.
La FANV è la più comune aritmia cardiaca sostenuta nella popolazione adulta: in
assenza di profilassi anticoagulante, il rischio di ictus ischemico è aumentato
da due a sette volte rispetto alla popolazione generale, ragione per cui
l'anticoagulazione a lungo termine è indicata per la maggior parte dei pazienti
con diagnosi confermata~\cite{esc2020afib}.

Il percorso decisionale della Nota segue quattro fasi sequenziali: conferma
elettrocardiografica della FANV; quantificazione del rischio tromboembolico
mediante il punteggio CHA\textsubscript{2}DS\textsubscript{2}-VASc; valutazione
del rischio emorragico con la scala HAS-BLED, che pesa fattori quali ipertensione, alterata funzione renale o epatica, pregresso ictus o sanguinamento ed età avanzata; selezione del farmaco e del
dosaggio appropriato in base alle comorbilità renali ed epatiche.
Il punteggio CHA\textsubscript{2}DS\textsubscript{2}-VASc assegna pesi a sette
condizioni cliniche: scompenso cardiaco ($+1$), ipertensione ($+1$), età
$\geq$75~anni ($+2$), diabete mellito ($+1$), pregresso ictus o attacco ischemico transitorio (TIA) ($+2$),
vasculopatia ($+1$) e sesso femminile ($+1$), per un massimo di 9 punti.
La soglia di rimborsabilità dipende dal sesso: un punteggio
CHA\textsubscript{2}DS\textsubscript{2}-VASc $\geq$2 nei maschi e $\geq$3 nelle
femmine è necessario per giustificare il trattamento rimborsabile.
Le controindicazioni assolute, ossia emorragia maggiore in atto, diatesi emorragica
congenita, gravidanza, ipersensibilità documentata al farmaco, protesi valvolare
meccanica o FANV valvolare (per i soli DOAC), producono un esito
\decision{NON\_RIMBORSABILE} immediato, verificato prima di qualsiasi altra
condizione.
La \notenova\ conta ventidue regole nel formato YAML del motore di regole ed è
l'unica delle quattro Note a richiedere l'intera gamma degli operatori logici
implementati, inclusa l'aritmetica d'intervallo per la gestione dei parametri
clinici parzialmente noti.

% ------------------------------------------------------------------
\section{Il problema della verifica della rimborsabilità}
\label{sec:problema-verifica}
% ------------------------------------------------------------------

La verifica della conformità di una prescrizione a una Nota AIFA è, nella pratica
corrente, affidata alla memoria e alla competenza del clinico senza alcun ausilio
sistematico.
L'errore prescrittivo non è un evento raro né privo di conseguenze: l'inappropriatezza
di una singola prescrizione si traduce in spesa farmaceutica non dovuta per il SSN
ma può anche determinare danno clinico quando una controindicazione assoluta,
come l'insufficienza renale grave nei pazienti trattati con determinati DOAC,
viene trascurata in favore dell'indicazione principale.
La letteratura documenta come la complessità cognitiva del processo prescrittivo,
amplificata dal numero e dalla densità normativa delle Note, sia associata a tassi
non trascurabili di non conformità con oneri tanto economici quanto di sicurezza
per il paziente~\cite{berner2007cdss}.

La difficoltà è strutturale.
Una Nota complessa come la \notenova\ richiede al medico di tenere simultaneamente
in mente la diagnosi di FANV, il calcolo di uno score a otto variabili con
aritmetica intera, le soglie differenziali per sesso, una tabella di dosaggio
dipendente dalla funzione renale e un insieme di controindicazioni assolute che
si applicano in modo disomogeneo ai quattro principi attivi.
Ogni parametro clinico non rilevato introduce un margine di incertezza che il
giudizio umano gestisce implicitamente, mentre un sistema formale richiede regole
esplicite per la propagazione del valore ignoto attraverso la catena decisionale.

L'automazione di questa verifica, mediante un sistema che integri la lettura delle
Note con i dati clinici del paziente e produca un esito motivato e tracciabile,
è la motivazione applicativa centrale di questa tesi.
La progettazione dell'architettura, la scelta del formalismo logico per la gestione
dell'incertezza e il protocollo di valutazione empirica sono descritti a partire
dal \Cref{ch:progettazione}.

È su questo terreno clinico e regolatorio che prendono forma le domande di ricerca
che guidano il lavoro. L'eterogeneità logica delle quattro Note (dalla semplice
congiunzione della \notezero\ allo score composito della \notenova) è ciò che rende
non banale la prima domanda, se cioè un componente simbolico possa replicare il
giudizio dell'annotatore umano sull'intera varietà dei casi (RQ1); la posta clinica
della verifica, in cui una controindicazione trascurata può tradursi in danno, è ciò
che impone alle domande successive di misurare non solo l'utilità del sistema, ma la
sua sicurezza, la fedeltà della spiegazione alle fonti e la tracciabilità dell'intero
percorso decisionale. Il capitolo seguente introduce gli strumenti tecnici con cui
queste domande verranno affrontate.
